\documentclass[10pt,a4paper]{amsart}
\usepackage[margin=19mm]{geometry}
\usepackage{amsmath,amssymb,mathtools}
\usepackage[scr=rsfs,cal=euler]{mathalfa}
\usepackage{xfrac}
\usepackage[unicode,hidelinks,bookmarks=false]{hyperref}
\newcommand{\ie}{i.e.\ }
\newcommand{\cf}{cf.\ }
\newcommand{\resp}{respectively\ }
\newcommand{\Subsection}{\subsection}
\providecommand{\nobreakdash}{\nobreak\hbox{-}}
\setlength{\emergencystretch}{2em}
\multlinegap0pt
\usepackage{amssymb}
\newtheorem{lemma}{Lemma}
\title{A single scale smooth Alpert trilinear characterization of the Fourier extension conjecture on $\mathbb{P}^{2}$}
\author{Cristian Rios, Eric T. Sawyer}
\date{Source: arXiv:2506.23376v5 (CC BY 4.0)}
\begin{document}
\maketitle

\noindent\emph{Source and licence.} A single scale smooth Alpert trilinear characterization of the Fourier extension conjecture on $\mathbb{P}^{2}$. Version arXiv:2506.23376v5 (CC BY 4.0), distributed by arXiv under the Creative Commons Attribution 4.0 International licence, http://creativecommons.org/licenses/by/4.0/, as recorded in the arXiv licence metadata for this exact version and shown on the abstract page https://arxiv.org/abs/2506.23376v5. Full licence notice: \texttt{licenses/arxiv-2506.23376-CC-BY-4.0.txt}.\par\smallskip

\noindent\textit{Editorial scope.} Counted unit: the article's lemma inv (source0.tex lines 863--872). The article's complete original proof of it is delivered with it (source0.tex lines 874--949, one \texttt{proof} environment of 76 lines). No mathematics is written, repaired or re-derived: the statement and the proof are the article's own lines, in the article's own notation, and no retained line is substituted or re-flowed.  Retained uncounted as the source-local context the proof needs: lines 531--534.  Nothing was omitted from any retained span, so the package carries no omission record.  No retained line contains a citation, so the wrapper carries no bibliography and no bibliography entry.  No retained line contains a figure environment, an image-inclusion command or drawing code, so no image is delivered and the package declares no asset dependency.  Editorial formatting only, in addition: an AMS \texttt{amsart} preamble that restates the article's own declarations and macros used by the retained lines, namely the environments \texttt{lemma} on separate counters, together with the standard packages those lines need.  The retained environments print in this document as Lemma 1 in order of appearance, whereas the article prints the same items with the numbering of its own full document.  The complete native source of the article as distributed in the arXiv e-print for this exact version is bundled unchanged as \texttt{sources/180703/source0.tex}.
\begin{equation}
\delta_{m,p}\left(  S_{\kappa,\eta}^{\mathcal{G}}\right)  ^{-1}=\left(
S_{\kappa,\eta}^{2^{m}\mathcal{G}}\right)  ^{-1}\delta_{m,p}\ .\label{claim}%
\end{equation}

\begin{lemma}
\label{proj inv}Let $1<q<p\leq\infty$. Then%
\begin{equation}
\delta_{m,p}\left(  \bigtriangleup_{I;\kappa}^{\eta,\mathcal{G}}f\right)
=\bigtriangleup_{2^{m}I;\kappa}^{\eta,2^{m}\mathcal{G}}\delta_{m,p}%
f,\ \ \ \ \ \text{if both }I,2^{m}I\in\mathcal{G}\left[  U\right]  \text{ and
}f\in L^{p}\left(  U\right)  ,\label{inv}%
\end{equation}

\end{lemma}

\begin{proof}
It suffices to prove the case $p<\infty$ since $L^{\infty}\left(  U\right)
\subset L^{p}\left(  U\right)  $. We have,%
\begin{equation}
\delta_{m,p}\left(  \bigtriangleup_{I;\kappa}^{\eta,\mathcal{G}}f\right)
\left(  x\right)  =\delta_{m,p}\left(  \left\langle \left(  S_{\kappa,\eta
}^{\mathcal{G}}\right)  ^{-1}f,h_{I;\kappa}\right\rangle h_{I;\kappa}^{\eta
}\right)  \left(  x\right)  =\left\langle \left(  S_{\kappa,\eta}%
^{\mathcal{G}}\right)  ^{-1}f,h_{I;\kappa}\right\rangle \delta_{m,p}%
h_{I;\kappa}^{\eta}\left(  x\right)  ,\label{one}%
\end{equation}
and%
\begin{equation}
\bigtriangleup_{2^{m}I;\kappa}^{\eta,2^{m}\mathcal{G}}\delta_{m,p}f\left(
x\right)  =\left\langle \left(  S_{\kappa,\eta}^{2^{m}\mathcal{G}}\right)
^{-1}\delta_{m,p}f,h_{2^{m}I;\kappa}\right\rangle h_{2^{m}I;\kappa}^{\eta
}\left(  x\right)  =\left\langle \delta_{m,p}\left(  S_{\kappa,\eta
}^{\mathcal{G}}\right)  ^{-1}f,h_{2^{m}I;\kappa}\right\rangle h_{2^{m}%
I;\kappa}^{\eta}\left(  x\right)  ,\label{two}%
\end{equation}
since $2^{m}I\in2^{m}\mathcal{G}$, and since $\left(  S_{\kappa,\eta}%
^{2^{m}\mathcal{G}}\right)  ^{-1}\delta_{m,p}=\delta_{m,p}\left(
S_{\kappa,\eta}^{\mathcal{G}}\right)  ^{-1}$ by (\ref{claim}).

We first compute the factor $\left\langle \delta_{m,p}\left(  S_{\kappa,\eta
}^{\mathcal{G}}\right)  ^{-1}f,h_{2^{m}I;\kappa}\right\rangle $ in (\ref{two})
in terms of $I$. We have%
\begin{align}
&  \ \ \ \ \ \ \ \ \ \ \ \ \ \ \ \ \ \ \ \ \ \ \ \left\langle \delta
_{m,p}\left(  S_{\kappa,\eta}^{\mathcal{G}}\right)  ^{-1}f,h_{2^{m}I;\kappa
}\right\rangle =\int\frac{1}{2^{\frac{2}{p}m}}\left(  S_{\kappa,\eta
}^{\mathcal{G}}\right)  ^{-1}f\left(  \frac{y}{2^{m}}\right)  h_{2^{m}%
I;\kappa}\left(  y\right)  dy\label{three}\\
&  =\int\frac{1}{2^{\frac{2}{p}m}}\left(  S_{\kappa,\eta}^{\mathcal{G}%
}\right)  ^{-1}f\left(  z\right)  h_{2^{m}I;\kappa}\left(  2^{m}z\right)
2^{2m}dz=2^{2m}\int\left(  S_{\kappa,\eta}^{\mathcal{G}}\right)  ^{-1}f\left(
z\right)  \frac{1}{2^{\frac{2}{p}m}}\frac{h_{2^{m}I;\kappa}\left(
2^{m}z\right)  }{h_{I;\kappa}\left(  z\right)  }h_{I;\kappa}\left(  z\right)
dz\nonumber\\
&  =2^{2m}\int\left(  S_{\kappa,\eta}^{\mathcal{G}}\right)  ^{-1}f\left(
z\right)  \frac{1}{2^{\frac{2}{p}m}}\sqrt{\frac{\left\vert I\right\vert
}{\left\vert 2^{m}I\right\vert }}h_{I;\kappa}\left(  z\right)  dz=\frac{2^{m}%
}{2^{\frac{2}{p}m}}\int\left(  S_{\kappa,\eta}^{\mathcal{G}}\right)
^{-1}f\left(  z\right)  h_{I;\kappa}\left(  z\right)  dz=2^{m\left(
1-\frac{2}{p}\right)  }\left\langle \left(  S_{\kappa,\eta}^{\mathcal{G}%
}\right)  ^{-1}f,h_{I;\kappa}\right\rangle .\nonumber
\end{align}
Next we compute the factor $\delta_{m,p}h_{I;\kappa}^{\eta}\left(  x\right)  $
in (\ref{one}) in terms of $2^{m}I$. We have%
\begin{equation}
\delta_{m,p}h_{I;\kappa}^{\eta}\left(  x\right)  =\frac{1}{2^{\frac{2}{p}m}%
}h_{I;\kappa}^{\eta}\left(  \frac{x}{2^{m}}\right)  =\frac{1}{2^{\frac{2}{p}%
m}}\frac{h_{I;\kappa}\left(  \frac{x}{2^{m}}\right)  }{h_{2^{m}I;\kappa
}\left(  x\right)  }h_{2^{m}I;\kappa}^{\eta}\left(  x\right)  =\frac
{1}{2^{\frac{2}{p}m}}\sqrt{\frac{\left\vert 2^{m}I\right\vert }{\left\vert
I\right\vert }}h_{2^{m}I;\kappa}^{\eta}\left(  x\right)  =2^{m\left(
1-\frac{2}{p}\right)  }h_{2^{m}I;\kappa}^{\eta}\left(  x\right)  .\label{four}%
\end{equation}
Using first (\ref{one}), then (\ref{four}), then (\ref{three}) and finally
(\ref{two}) gives (\ref{inv}),%
\begin{align*}
\delta_{m,p}\left(  \bigtriangleup_{I;\kappa}^{\eta,\mathcal{G}}f\right)
\left(  x\right)   &  =\left\langle \left(  S_{\kappa,\eta}^{\mathcal{G}%
}\right)  ^{-1}f,h_{I;\kappa}\right\rangle \delta_{m,p}h_{I;\kappa}^{\eta
}\left(  x\right) \\
&  =\left\langle \left(  S_{\kappa,\eta}^{\mathcal{G}}\right)  ^{-1}%
f,h_{I;\kappa}\right\rangle 2^{m\left(  1-\frac{2}{p}\right)  }h_{2^{m}%
I;\kappa}^{\eta}\left(  x\right) \\
&  =\left\langle \delta_{m,p}\left(  S_{\kappa,\eta}^{\mathcal{G}}\right)
^{-1}f,h_{2^{m}I;\kappa}\right\rangle h_{2^{m}I;\kappa}^{\eta}\left(  x\right)
\\
&  =\bigtriangleup_{2^{m}I;\kappa}^{\eta,2^{m}\mathcal{G}}\delta_{m,p}f\left(
x\right)  .
\end{align*}

\end{proof}

\end{document}
